\subsection{Bernoulli units and the rank-four arithmetic data}
\label{sec:selmer-dyadic-twist-vanishing}

The rank-four construction uses an auxiliary odd character of large
dyadic conductor to separate the residual characters, even when the
original character is trivial on the coefficient-prime decomposition
group. The following elementary argument gives the full Bernoulli
unit threshold required in that construction, with the odd conductor
factor retained.

\begin{theorem}[Uniform Bernoulli units for dyadic twists]
\label{thm:selmer-dyadic-twist-unit}
Let $\ell$ be odd, $d\ge1$, $Q=\ell^d$, and
$h=v_2(Q-1)\ge2$. Let $\varphi$ be a primitive Dirichlet
character of odd conductor $f$ and order dividing $Q-1$,
and let $\xi$ be primitive of conductor $M=2^n$.
Assume
\[
       \varphi(-1)\xi(-1)=-1,\qquad f4^h\le M.        \tag{BU1}
\]
Then $B_{1,\varphi\xi}$ is an $\ell$-adic unit at every
prime of its character-value field above $\ell$.
This includes $\ell\mid f$ and is the full source
\texttt{LemmaX} assertion.
\end{theorem}

\begin{proof}
Work first in characteristic zero, with the chosen
$\ell$-adic character values. For $u\in\mathbf Z/M\mathbf Z$
let $\tilde u$ be its representative in $[0,M)$, and put
\[
 A(u)=\sum_{\substack{0\le a<fM\\a\equiv u\bmod M}} a\varphi(a),
 \qquad S=\sum_{a=0}^{f-1}\varphi(a),\qquad
 g(u)=\frac{A(u)-A(0)}{fM}.
\]
The sum of $\varphi(a)$ over the congruence class used
in $A(u)$ is $S$, independently of $u$, since $f$ and $M$
are coprime. Shifting the summation by $ft$ modulo $fM$
gives, for $0\le t<M$,
\[
 g(u)-g(u-ft)=\frac{tS}{M}-
       \sum_{\substack{0\le a<ft\\a\equiv u\bmod M}}\varphi(a).
                                                               \tag{BU2}
\]
Indeed each shifted integer increases by $ft$ and decreases
by $fM$ precisely when it wraps around the interval.
The character is unchanged by this shift. The formula
uses the same integer representatives in both sums.

For each $u$ choose $t$ with $ft\equiv u\bmod M$.
As $g(0)=0$, (BU2) proves that $g(u)$ is integral and
its reduction belongs to $\mathbf F_Q$: the reduced
values of $\varphi$ lie there and $M$ is a unit.
This step avoids division by the odd conductor in the
residue field, including when $\ell\mid f$.
The character $\xi$ is nontrivial by (BU1), which implies
$n\ge4$. Hence $\sum_u\xi(u)=0$. Grouping the Bernoulli
sum by residue classes modulo $M$ gives the exact identity
\[
              B_{1,\varphi\xi}=\sum_u \xi(u)g(u).               \tag{BU3}
\]
In particular this Bernoulli number is integral.

From now on write the reduced values of $\xi,\varphi,g,S$
without bars, in an algebraic closure of $\mathbf F_\ell$.
The character $\xi$ is still primitive, since reduction
preserves roots of unity of power-of-two order.
Suppose (BU3) reduces to zero. Put $m=n-h$ and
$s=f2^h$, so that $m\ge h\ge2$ and $s\le2^m<M$.
For every unit $y$ modulo $M$, changing variables gives
\[
                         \sum_x \xi(x)g(xy)=0.                 \tag{BU4}
\]
We extract part of this identity by the Frobenius action.
The order of $Q$ modulo $M$ is $2^{n-h}$, and its powers
are exactly the residues congruent to one modulo $2^h$.
This follows from
$v_2(Q^{2^j}-1)=h+j$, by successive squaring since
$Q\equiv1\bmod4$. Raise (BU4) to these powers of Frobenius.
Because $g$ takes values in $\mathbf F_Q$, this gives
$\sum_x\xi(x)^{1+2^h j}g(xy)=0$ for
$0\le j<2^{n-h}$. Summing and dividing by $2^{n-h}$,
which is a unit, gives by a geometric sum
\[
 \sum_{x\in H}\xi(x)g(xy)=0,\qquad
 H=\{x\in(\mathbf Z/M\mathbf Z)^\times:
                                  \xi(x)^{2^h}=1\}.            \tag{BU5}
\]
The cyclic group $1+4\mathbf Z/M\mathbf Z$ is generated
by $5$, and primitivity of $\xi$ makes it faithful on
that group. It follows that
\[
 H=H_+\sqcup H_-,\qquad
 H_\pm=\{x:x\equiv\pm1\bmod 2^m\}.
\]
Each of the two sets has character sum zero: multiplication
by $1+M/2\in H_+$ preserves both sets and changes the sign
of $\xi$. This also covers $m=2$.

Apply (BU5) at $y=1$ and $y=1-s$, both units. Write
$C=2^h S/M$. For $x\in H_+$ we have $xs=s\bmod M$;
the only member with representative less than $s$ is
$1$. Thus (BU2) gives
\[
             g(x)-g(x(1-s))=C-\mathbf 1_{x=1}.
\]
For $x\in H_-$ we have $xs=-s\bmod M$. The only member
for which $\widetilde{x+s}<s$ is $x=-1$, with
$\widetilde{x+s}=s-1$. Applying (BU2) at $x+s$ gives
\[
       g(x)-g(x(1-s))=-C+\mathbf 1_{x=-1}\varphi(s-1).
\]
These assertions follow directly from $s\le2^m$ and
the congruences defining $H_\pm$; they include equality
in (BU1). Since $s$ is a multiple of $f$,
$\xi(-1)\varphi(s-1)=-1$. Subtracting the two instances
of (BU5), the terms involving $C$ cancel by the character
sums just established, while the exceptional terms give
$-1-1=-2$. This is nonzero in characteristic $\ell$,
a contradiction. Thus (BU3) has nonzero residue.
\end{proof}

\begin{corollary}[Unramified cocycles in the original character]
\label{cor:selmer-dyadic-twist-cocycle}
Retain (BU1), and let $\gamma$ be the reduction of
$\varphi\xi$ or its inverse in a field $k$ of
characteristic $\ell$. A continuous cocycle
$z:G_{\mathbf Q}\to k(\gamma)$, with the discrete topology
on $k$, is a coboundary if for every finite inertia group
$I$ it vanishes on a subgroup of $I\cap\ker\gamma$
of finite index prime to $\ell$.

This proves the original \texttt{CorX} with its locally
constant cocycle and relative-index conditions.
No assumption that $\ell$ is prime to the degree of
the full cyclotomic field is required.
\end{corollary}

\begin{proof}
We first give the characterwise Stickelberger argument
with its necessary hypotheses. Write $\chi=\varphi\xi$
or $\chi=(\varphi\xi)^{-1}$ so that $\gamma=\bar\chi$.
Set $F=\mathbf Q(\mu_{fM})$, $G=\operatorname{Gal}(F/\mathbf Q)$.
Restriction of $z$ to $G_F$ is a homomorphism, with
finite image since its domain is compact and its
target is discrete. It vanishes on inertia in $G_F$:
if $I_F=I\cap G_F$ and the given subgroup is
$J\subset I\cap\ker\gamma$, the index
$[I_F:I_F\cap J]$ divides $[I\cap\ker\gamma:J]$ because
$I_F$ is normal in $I\cap\ker\gamma$.
The image of $I_F$ in the additive group of $k$ thus
has both $\ell$-power order and order prime to $\ell$.
It is zero. If the original relative-index condition
is specified by an arbitrary ambient subgroup, first
intersect it with $I\cap\ker\gamma$.
Class field theory identifies the restricted homomorphism
with an additive map $u:\operatorname{Cl}(F)\to k$.
There are no real places of $F$. The cocycle relation
and reciprocity give
\[
                  u(\sigma_b x)=\bar\chi(b)u(x)
       \quad ((b,fM)=1).
\]

For $D=fM$, the classical Stickelberger element is
$\Theta=D^{-1}\sum_{(a,D)=1,\ 1\le a<D}a\sigma_a^{-1}$.
For any integer $b$ prime to $D$, $(b-\sigma_b)\Theta$
is integral (subtract the least positive residues of
$ba$ modulo $D$). Stickelberger's theorem says that
this integral group-ring element annihilates
$\operatorname{Cl}(F)$ \cite[Theorem 4.5.6]{dyadic-bernoulli-sharifi}.
Its scalar under $\chi$ is
$(b-\chi(b))B_{1,\chi^{-1}}$.
Theorem~\ref{thm:selmer-dyadic-twist-unit} gives the
unit assertion for $\chi^{-1}$, since (BU1) is preserved
by inversion. There is $b$ with $b-\bar\chi(b)\ne0$
in $k$: if $\ell\nmid D$ take $b=\ell$; otherwise
$\bar\chi$ differs from the residual cyclotomic
character modulo $D$, because its dyadic component
is primitive of conductor $2^n\ge16$ while the
cyclotomic character has conductor $\ell$.
Thus the displayed integral annihilator acts by a
nonzero scalar on $u$, and $u=0$.

Now $z$ descends to a cocycle of the finite abelian
group $G$. Choose $g_0$ with $\gamma(g_0)\ne1$,
which exists by the high dyadic conductor.
Commutativity gives
\[
 (\gamma(g_0)-1)z(g)=(\gamma(g)-1)z(g_0).
\]
Division by the nonzero scalar on the left proves
that $z$ is a coboundary. In particular this step
does not require averaging over the full group $G$,
whose order can be divisible by $\ell$.
\end{proof}

The same Stickelberger argument proves the underlying
unramified-cocycle implication whenever the inverse
Bernoulli value at the actual modulus is a unit, the
finite character has order prime to $\ell$ and its
reduction is distinct from both the trivial and
residual cyclotomic characters. These are exactly the
nontriviality and cyclotomic-exclusion hypotheses in the
source \texttt{UnramifiedCocycleIsCoboundaryOfBernoulliUnit}.
The conductor hypothesis of the preceding theorem is
only needed to supply that unit.

\begin{theorem}[The original RH rank-four arithmetic datum]
\label{thm:rh-symplectic-arithmetic-data}
Retain the actual ordinary RH host with modulus
$N=2^a qd$, a nontrivial character $\psi$ of its
Galois group, an odd coefficient prime $\ell\nmid N$,
and $e\ge2$ with $\psi(-1)=(-1)^e$.
Write $\psi=\psi_2\psi_{\rm odd}$ with dyadic conductor
$2^b$ and $\psi_{\rm odd}^2=1$. Put
\[
 d_0=\operatorname{ord}_4(\ell),\quad
 h=v_2(\ell^{d_0}-1),\quad
 n=\left\lceil\log_2\max\{\ell N4^h,
                                      2^{b+6}\ell^3\}\right\rceil.
                                                               \tag{BU6}
\]
These are exactly the original RH specialization of
\texttt{f0}, \texttt{n0} and \texttt{Host.nII}.
After the permitted finite extension of the coefficient
DVR containing the required roots of unity, there is
an original \texttt{Raw} datum satisfying every field of
\texttt{Good} and the original \texttt{OfChar} for Case II
(with $p_1=m_0=1$). Its dyadic modulus satisfies
\[
                            2^n\le128N\ell^5.                  \tag{BU7}
\]
In particular this constructs the actual rank-four
arithmetic datum for every remaining character of (DC11).
It does not assert automorphy or integral rigidity of
its deformations.
\end{theorem}

\begin{proof}
The finite character $\psi$ has power-of-two order;
its odd component is quadratic or trivial. Consequently
the source's prime-to-$\ell$ part is $\psi$ itself,
its $\ell$-component is trivial, and its \texttt{th1}
is $\psi_{\rm odd}$. This proves the formulas for
$d_0,h,n$ in (BU6). Choose an odd primitive character
$\xi$ of conductor $2^n$, with its values in a finite
coefficient extension. Such a character is obtained
by assigning a primitive $2^{n-2}$-th root to $5$ and
$-1$ to $-1$. The added roots have order prime to
$\ell$, so adjoining them can be done unramified.
The length normalization is unchanged under this
extension by the integral coefficient comparison
already used for the RH character groups.

In the source's notation take
\[
 \Psi=\varepsilon^e\psi,\quad
 \omega_T=\text{Teichm\"uller lift of }\bar\varepsilon,\quad
 \alpha=\varepsilon\omega_T^{-1}\xi,\quad
 \beta=\Psi\alpha^{-1},\quad \tau=\alpha\Psi .
\]
Here $\varepsilon$ is the $\ell$-adic cyclotomic character.
These are precisely \texttt{Psi}, \texttt{om},
\texttt{al}, \texttt{be}, \texttt{th}, respectively.
Take the original set $S$ of prime divisors of
$2\ell N$, choose absolute places over the rational
primes, and take the original complex conjugation $c$.
The base diagonal is $(\tau,\Psi,\alpha,1)$, with
similitude character $\tau$ and weights $(e+1,e,1,0)$.
The parity condition gives $\Psi(c)=1$,
$\alpha(c)=\beta(c)=\tau(c)=-1$.
The orders and ramification of the finite characters
are as required. The source's exceptional set
\texttt{Sb} is empty because $\psi$ has order prime
to $\ell$; its \texttt{T} is $\{\ell\}$.
In particular $\xi$ is unramified at every place of
\texttt{T}.

Choose an arithmetic Frobenius lift $s_\ell$ at $\ell$
with $\varepsilon(s_\ell)=1$. It exists because the
cyclotomic $\ell$-power extension of $\mathbf Q_\ell$
is totally ramified, and the cyclotomic character on
inertia is onto $\mathbf Z_\ell^\times$. Then
$\Psi(s_\ell)=\psi(\ell)$ and $\xi(s_\ell)=\xi(\ell)$.
The valuation identity
\[
 v_2(\ell^{2^j}-1)=v_2(\ell^2-1)+j-1
                   \quad(j\ge1)
\]
and (BU6) imply
$n>v_2(\ell^2-1)+b+1$: indeed
$v_2(\ell^2-1)<2\log_2\ell$ and
$n\ge b+6+3\log_2\ell$.
Thus $\ell^{2^{b+2}}\not\equiv1\bmod2^n$.
This power lies in $1+4\mathbf Z/2^n\mathbf Z$;
faithfulness of the primitive character there gives
$\xi(\ell)^{2^{b+2}}\ne1$.
The order of $\psi$ divides $2^{b+2}$, including
$b=0$. All these roots have order prime to $\ell$,
so their orders are preserved by reduction.
Consequently
\[
       \bar\xi(s_\ell)\bar\Psi(s_\ell)^i\ne1
                         \quad(i\in\mathbf Z).                 \tag{BU8}
\]
This proves the source \texttt{condD} for every
integer $i$, not merely the finitely many matrix weights.

The three reduced characters
$\bar\alpha,\bar\beta,\bar\tau$ are the reductions
of the finite characters
\[
             \xi,\qquad \omega_T^e\psi\xi^{-1},
                           \qquad \omega_T^e\psi\xi.
\]
All are odd. Their dyadic components have exact conductor
$2^n$, since $n>b$ and multiplication by a lower-conductor
character does not change primitivity there. Their
primitive odd components have conductors dividing
$\ell qd$ and orders dividing $\ell^{d_0}-1$.
This also holds for their inverses.
Thus (BU1) holds for each, by (BU6).
Suppose a cocycle satisfies the source's \texttt{XVan}
premises for one of these six characters.
Away from $\ell$ it vanishes on inertia intersected
with the kernel of the full residual diagonal; that
kernel has prime-to-$\ell$ index in the inertia group
and is contained in the character kernel in question.
At $\ell$ its local class is a coboundary, hence the
cocycle is zero on inertia intersected with the
character kernel. The same local conclusion holds at
the conjugate places by the cocycle identity.
Corollary~\ref{cor:selmer-dyadic-twist-cocycle} now
gives the required global coboundary.
This proves all six original \texttt{xvan} assertions.

For the original element $\mathrm{ph}$, the sequence
$s_\ell^{\ell^{j!}}$ converges in the profinite group:
in every finite quotient the $\ell$-primary component
of $s_\ell$ is eventually killed and its prime-to-$\ell$
component is eventually unchanged, because
$\ell^{j!}\equiv1$ modulo its order for large $j$.
These compatible limits define the desired element.
For the second mark choose $g_0$ outside the kernels of
both $\bar\alpha$ and $\bar\Psi$. These are proper
subgroups: the first character has high dyadic
conductor, and the second is nontrivial since $\psi$
is a nontrivial prime-to-$\ell$ character unramified
at $\ell$ (the residual cyclotomic factor cannot cancel
its ramification at a different finite prime).
The union of two proper subgroups is not a group,
so such $g_0$ exists. Finally choose $s_0$ in inertia
at $\ell$ with $\varepsilon(s_0)=1+\ell$.
It kills both $\psi$ and $\omega_T$ and has infinite
cyclotomic order. This is the original
\texttt{decompSigma} condition.
All other fields of \texttt{Raw.Good} are the
specified absolute decomposition and inertia groups,
the ramification and parity conditions checked above,
and the exact finite set of places. The character
$\xi$ is the required \texttt{OfChar} character
with the odd component of conductor one. This verifies
the actual datum without changing its modulus.

Lastly $4^h\le(\ell^{d_0}-1)^2<\ell^4$ and $2^b\le N$.
Since the least power of two above a positive integer
is less than twice it, (BU6) gives (BU7).
The exponent and its constants are independent of
the original fields and of $e$.
\end{proof}

This completes the original arithmetic input to the
rank-four Case II argument for the remaining RH family,
and the Bernoulli-unit and unramified-cocycle inputs for
both RH and GRH with their exact original character
hypotheses. The remaining rank-four tasks concern
the actual local components, the characteristic-prime
finiteness, automorphy and rigidity of the relevant
points, and their integral length comparison. None
follows just from (BU7) or from the six residual vanishings.
